%% v4/sec/methodology.tex — §5 Benchmark Methodology (TODAES long-form)
%%
%% Provenance: verbatim mechanical carry of paper/sec/methodology.tex
%% (2026-05-16) under Phase-1 task 91240545. v3 source is frozen.
%%
%% PHASE-2 expansion targets (from outline_v4.md §5 plan, ~2.0-2.5pp):
%%   - 5.1 Workload — expand circuit-selection rationale (why
%%     LC_VCO_20G as headline + OpAmp as transferability check);
%%     surface the foundry-NDA constraint as design-driver
%%   - 5.2 Spec contract — promote YAML metric contract to first
%%     subsection with full schema walkthrough + alternative
%%     considered (gymspec / TestbenchCAD-DSL) + chosen-mechanism
%%     rationale
%%   - 5.3 Protocol — inline JSON tool contract (currently in
%%     supplementary); per-design-choice rationale for stuck-streak
%%     guard / iteration budget cap / replicate key
%%   - 5.4 Seed handling — expand to address reviewer concern
%%     "is replicate key really not a sampling input?"
%%   - 5.5 Repetition — un-trim per Q-mig-2: justify $11\times\geq 3$
%%     floor as statistical power floor not arbitrary
%%   - 5.6 Three-axis audit — canonical-test walkthrough with
%%     assertion-by-assertion explanation (cross-ref §4.5 stub)
%%   - 5.7 Metric definitions — expand Anthropic / Gemini
%%     reasoning-token un-capture limitation
%%   - 5.8 Threats to validity — promote to dedicated subsection
%%     with per-threat mitigation paragraph

\section{Benchmark Methodology}
\label{sec:methodology}

\subsection{Workload}
\label{sec:method-workload}

We evaluate on \textbf{LC\_VCO\_20GHz}, a 20\,GHz inductively-loaded
differential VCO laid out on a real industrial-node CMOS PDK
(specific node withheld under NDA; reported in supplementary
material with metric scales normalised to preserve compliance).
An additional OpAmp testbench is reported in revision/appendix if
available before deadline.
%% PHASE-2: expand circuit-selection rationale + foundry-NDA constraint as design-driver.

\subsection{Spec contract}
\label{sec:method-spec}

The spec is a YAML metric contract (signals / windows / metrics)
identical across all eleven LLMs; the seven concrete metrics for
LC\_VCO\_20G are the machine-readable contract the evaluator
consumes (Table~\ref{tab:spec}, taken from
\filename{spec.md:61--88}). The evaluator supports only the simple
stats \texttt{\{mean, min, max, ptp, rms, mean\_abs, freq\_Hz,
duty\_pct\}} plus the compound forms
\texttt{\{ratio, t\_cross\_frac\}}
(\filename{src/spec\_evaluator.py:104-106}); phase-noise / K\_VCO
are out of scope for v3. A run \textbf{passes} iff every metric is
within its \texttt{pass} range; failures classify by which metric
fell out of range (\texttt{MEASURED\_OUT\_OF\_RANGE}) versus the
evaluator declaring the measurement infeasible
(\texttt{UNMEASURABLE}).
%% PHASE-2: promote YAML schema walkthrough + alternative considered (gymspec / TestbenchCAD-DSL).

\begin{table}[t]
  \centering
  \caption{LC\_VCO\_20G spec contract (\filename{spec.md:61--88}).
    A run passes iff every metric is inside its \texttt{pass}
    range; \texttt{[null,\ X]} denotes upper-only,
    \texttt{[X,\ null]} lower-only. The \texttt{sanity} range is
    the topology-feasibility window outside which the run is
    classified as a topology-induced FAIL (\S\ref{sec:failures}).}
  \label{tab:spec}
  \footnotesize
  \setlength{\tabcolsep}{3.5pt}
  \begin{tabular}{l l l l}
    \toprule
    \textbf{Metric} & \textbf{Window} & \textbf{Pass} & \textbf{Sanity} \\
    \midrule
    \texttt{f\_osc\_GHz}      & full         & $[19.5,\ 20.5]$    & $[0.1,\ 100]$ \\
    \texttt{V\_diff\_pp\_V}   & late         & $[0.40,\ \text{null}]$ & $[0,\ 8]$ \\
    \texttt{V\_cm\_V}         & late         & $[0.70,\ 0.81]$    & $[0,\ 1.6]$ \\
    \texttt{duty\_cycle\_pct} & late         & $[48,\ 52]$        & $[0,\ 100]$ \\
    \texttt{amp\_hold\_ratio} & late/early   & $[0.95,\ \text{null}]$ & $[0,\ 3]$ \\
    \texttt{t\_startup\_ns}   & startup      & $[\text{null},\ 10]$ & $[0,\ 200]$ \\
    \texttt{I\_core\_uA}      & late         & $[\text{null},\ 800]$ & $[0,\ 10000]$ \\
    \bottomrule
  \end{tabular}
\end{table}

\subsection{Protocol}
\label{sec:method-protocol}

For each (LLM, run-seed) tuple the agent runs the same Safe-Bridge
configuration, the same prompt template, and the same JSON tool
contract; only the \texttt{-{}-llm} flag differs across cells of the
matrix. Each run iterates until either (i)~the spec evaluator
reports first-pass, or (ii)~a stuck-streak guard (no design-vector
change across 3 successive iterations) aborts, or (iii)~an
iteration budget cap of $N$ (resolved per pilot) is reached. The eleven
checkpoints span seven vendor families---Anthropic Opus 4.7 /
Sonnet 4.6 / Haiku 4.5 (US), OpenAI GPT-5.5 / GPT-5.4-mini (US),
Google Gemini 2.x (US), Moonshot Kimi K2.5 (China),
MiniMax M2.7 (China), Xiaomi MiMo v2.5-pro (China), DeepSeek
V4-pro / V4-flash (China)---drawn from a six-US + five-China split
(\S\ref{sec:platform-threat}).
%% PHASE-2: inline JSON tool contract + per-design-choice rationale (stuck-streak / budget cap / replicate key).

\subsection{Seed handling and labelling}
\label{sec:method-seed}

Each run uses a replicate key from $\{1, 2, 3\}$
(\filename{scripts/run\_benchmark.py:94});
\filename{build\_subprocess\_cmd}
(\filename{scripts/run\_benchmark.py:187-217}) does not pass the
seed to \filename{run\_agent.py} or any API request. The spec
contract is byte-for-byte identical across LLMs, so seed acts
purely as a replicate label, not as a prompt or sampling-API
input.
%% PHASE-2: expand to address reviewer concern "is replicate key really not a sampling input?"

\subsection{Repetition and variance}
\label{sec:method-rep}

Each LLM is run for at least \textbf{three independent seeds} per
cell of the matrix (more for LLMs identified as noisy in pilot).
Reported quantities are \textbf{median + interquartile range
(IQR)}, not best-of; single-run reports are flagged as such in
\S\ref{sec:results}. The $11\times\geq 3$ floor is 33 runs minimum.
%% PHASE-2: un-trim per Q-mig-2 — justify 11x>=3 as statistical-power floor (canonical un-trim source: paper/sec5_results.md).

\subsection{Three-axis scrub-audit methodology}
\label{sec:method-audit}

Every scrub-coverage regression test asserts on three independent
axes simultaneously: (a)~\textbf{records-present} (expected log
record emitted at expected level); (b)~\textbf{forbidden-tokens-%
absent} (every PDK-token set member absent from emitted text);
(c)~\textbf{scrub-marker-present}
(\texttt{<redacted>} / \texttt{<path>} marker present, confirming
the scrubber ran rather than the leak path being unreachable in
the test fixture). A single negative assertion (``token absent'')
false-positives when the log line is never emitted; the three-axis
pattern catches that failure mode. Canonical implementation:
\filename{test\_debug\_log\_scrubs\_thinking\_pdk\_tokens}
(\filename{tests/test\_llm\_client.py:147--190}); reasoning-trace
scrub call sites at \filename{llm\_client.py}
L339, L431, L533, L636, L742, L810, L818
(Kimi / MiniMax / OpenAI / MiMo / DeepSeek / Ollama-debug /
Ollama-fallback respectively).
%% PHASE-2: assertion-by-assertion canonical-test walkthrough; cross-ref §4.5 stub.

\subsection{Metric definitions}
\label{sec:method-metrics}

\begin{table}[t]
  \centering
  \caption{Headline metric definitions (Table~2).}
  \label{tab:metrics}
  \footnotesize
  \begin{tabularx}{\linewidth}{l X}
    \toprule
    \textbf{Metric} & \textbf{Definition} \\
    \midrule
    Success rate & Fraction of runs reaching all-metric-pass
      ($\in[0,1]$); reported as median across seeds. \\
    Iter-to-converge & Iteration count of first all-metric-pass
      ($\infty$ if never). \\
    Input tokens & Sum across all turns; per-vendor counter as
      documented in each provider's SDK. \\
    Output tokens & Sum across turns; for OpenAI excludes
      \texttt{reasoning\_tokens}; for Kimi / MiniMax / MiMo /
      DeepSeek excludes \texttt{reasoning\_content} (Anthropic /
      Gemini accounting is input + output + total only in v3, see
      \S\ref{sec:results-premium} caveat). \\
    Reasoning tokens & Broken out separately for vendors whose API
      surfaces them on the current \filename{\_normalize\_usage}
      schema (\filename{src/llm\_client.py:21-80}); Anthropic /
      Gemini reasoning-token usage is \textbf{not currently
      tracked} and is flagged as a v3 limitation in
      \S\ref{sec:results-premium} / \S\ref{sec:conclusion}. \\
    USD per run & Tokens $\times$ public pricing as of
      \textbf{2026-05-01}; pricing table in appendix; reasoning
      tokens billed at the documented per-vendor rate for vendors
      with tracked reasoning-token counts. \\
    Wall-clock / iter & Median over completed iterations; excludes
      rate-limit / 429 retry waits, separately reported. \\
    \bottomrule
  \end{tabularx}
\end{table}
%% PHASE-2: expand Anthropic / Gemini reasoning-token un-capture limitation as standalone callout.

\subsection{Threats to validity}
\label{sec:method-threats}

We enumerate the four primary validity threats reviewers most
often raise on benchmark papers, each with our intended response in
\S\ref{sec:results}--\S\ref{sec:discussion}:
%% PHASE-2: promote each threat to its own paragraph with mitigation argument.

\begin{enumerate}
  \item \textbf{Single-circuit scope.} Generalisation to other
    analog topologies is unproven. We state this up-front and use
    the ``OpAmp in revision/appendix'' hook
    (\S\ref{sec:method-workload}) to flag the open problem without
    promising a result we have not yet measured.

  \item \textbf{Pricing drift.} Public per-token pricing changes;
    we pin all USD figures to a single 2026-05-01 pricing snapshot,
    with the raw token counts also reported so re-pricing under
    future tariffs is straightforward.

  \item \textbf{Rate-limit / 429 noise.} Inflates wall-clock for
    cost-floor LLMs (Haiku 4.5, GPT-5.4-mini, DeepSeek V4-flash,
    MiMo v2.5-pro) asymmetrically. We separate rate-limit wait time
    from compute wall-clock and report both.

  \item \textbf{Reasoning-token accounting per-vendor differences.}
    Kimi, MiniMax, OpenAI, MiMo, and DeepSeek all return an
    OpenAI-compatible usage shape;
    \filename{\_normalize\_usage} reads the
    \texttt{completion\_tokens\_details.reasoning\_tokens}
    sub-field and surfaces it as the dedicated
    \texttt{reasoning\_tokens} counter
    (\filename{src/llm\_client.py:47-59}). The associated
    \texttt{reasoning\_content} \emph{text} (assistant reasoning
    trace, distinct from the usage counter) is the
    scrubbed-fallback path handled separately at the per-vendor
    reply sites (per-vendor reply sites enumerated in
    \S\ref{sec:method-audit}). Anthropic's thinking-token accounting
    (\texttt{cache\_creation\_input\_tokens} /
    \texttt{cache\_read\_input\_tokens}) and Gemini's
    \texttt{thinking}-block usage are \textbf{not currently
    captured} by \filename{\_normalize\_usage}
    (\filename{src/llm\_client.py:42-63}); the Anthropic and
    Gemini rows of the \S\ref{sec:results-premium} premium table
    are therefore marked unavailable, with a schema extension
    flagged in \S\ref{sec:conclusion}.

\end{enumerate}

\noindent Additional threats (provider-side caching,
cohort-selection bias, prompt-template optimality,
sampling-parameter standardisation) are discussed in the
supplementary appendix.
